\documentclass[12pt]{article}

% --------------------
% PACOTES ESSENCIAIS
% --------------------
\usepackage[utf8]{inputenc}     % Codificação UTF-8
\usepackage[T1]{fontenc}        % Fonte com acentuação
\usepackage[brazil]{babel}      % Idioma português
\usepackage{amsmath, amssymb}   % Símbolos matemáticos
\usepackage{graphicx}           % Inserção de imagens
\usepackage{geometry}           % Margens personalizadas
\usepackage{fancyhdr}           % Cabeçalhos e rodapés
\usepackage{enumitem}           % Listas personalizadas
%\usepackage{physics}            % Notação física e matemática
\usepackage{hyperref}           % Links no PDF
\usepackage{multicol}           % Colunas (opcional)

% --------------------
% CONFIGURAÇÕES
% --------------------
\geometry{a4paper, margin=2.5cm}

\pagestyle{fancy}
\fancyhf{}
\rhead{Método das Diferenças Finitas}
\lhead{Caio Magno Mendonça Nicácio}
\cfoot{\thepage}

\setlength{\parskip}{0.5em}
\setlength{\parindent}{0pt}

% --------------------
% INÍCIO DO DOCUMENTO
% --------------------
\begin{document}
	
	\begin{center}
		\Large\textbf{Método das Diferenças Finitas} \\
		\normalsize
		%Psicopatologia Laboral \\
		Data: \today
	\end{center}
	
	\hrule
	\vspace{1em}
	
	\tableofcontents
	\newpage
	
	\section{Polinômio de Taylor}
	
	Seja $f \in C^n$, $f:\mathbb{R} \to \mathbb{R}$, uma função diferenciável em $x_0$, $x,x_0 \in D_f$. O polinômio que melhor aproxima $f$ localmente, para todo $x$ na vizinhança de $x_0$, é descrito pela Eq. (\ref{eq:1}).
	\begin{equation}
		P_n(x)=f(x_0)+\sum_{n=0}^{k}\frac{f^{(n)}}{n!}(x-x_0)^n
		\label{eq:1}
	\end{equation}
	Seja $f \in C^2$, $f:\mathbb{R} \to \mathbb{R}$, uma função diferenciável no intervalo $I$, $x,x_0 \in I$.
	$$P_2(x)=f(x_0)+f'(x_0)(x-x_0)+\frac{f''(x_0)}{2}(x-x_0)^2$$
	O polinômio de Taylor de ordem 2 é a única função que goza da propriedade do erro $E_2(x)$ tender a zero mais rapidamente que $(x-x_0)^2$, tal que
	$$\lim_{x \to x_0} \frac{E_2(x)}{(x-x_0)^2}=0\text{,} \quad E_2(x)=o((x-x_0)^2)\text{.}$$
	Para valores de $\Delta x$ suficientemente pequenos, pode-se assumir que o termo quadrático é desprezível. Logo,
	$$P(x)=f(x_0)+f'(x_0)(x-x_0)\text{,} \quad f'(x_0) \neq 0\text{,}$$
	é a melhor aproximação local de $f$ em volta de $x_0$, quando $x \to x_0$. 
	
	Por definição, $P(x)=f(x)$, $\forall x \to x_0$. Portanto,
	$$f(x)\approx f(x_0)+f'(x_0)(x-x_0)\text{.}$$
	Isolando o termo diferencial $f'(x_0)$,
	$$\left.\frac{df}{dx}\right|_{x_0}\approx \frac{f(x)-f(x_0)}{x-x_0}$$
	para todo $x$ suficientemente próximo a $x_0$.
	
	\subsection{Aproximação Local das Derivadas de Ordem Superior}
	Seja $f \in C^{(n+1)}$, então a $n$-ésima derivada do polinômio  de Taylor de ordem $n$, $P_n^{(n)}(x)$, é uma boa aproximação local de $f^{(n)}(x)$ em volta de $x_0$. Por construção, portanto,
	$$P_n^{(k)}(x_0)=f^{(k)}(x_0)\text{,}\quad  k=0,1,...,n\text{.}$$
	Logo, para todo $x$ suficientemente próximo a $x_0$, $P_n^{(k)}(x)\approx f^{(k)}(x)$.
	O erro na aproximação da derivada $f^{(n)}(x)$ por $P_n^{(n)}(x)$ é:
	$$f^{(k)}(x)-P_n^{(k)}(x)=E_n^{(k)}(x)\text{.}$$
	Seja $f \in C^{(n+1)}$, então, pela forma de Lagrange do erro:
	$$E_n(x)=\frac{f^{(k+1)}(\xi)}{(k+1)!}(x-x_0)^{n+1}\text{,}\quad \xi \in (x_0,x)$$
	Derivando $n$ vezes,
	$$E_n^{(k)}=\frac{d^k}{dx^k} \left[\frac{f^{(k+1)}(\xi)}{(k+1)!}(x-x_0)^{k+1}\right]$$
	$$\frac{d^k}{dx^k}(x-x_0)^{k+1}=\frac{(k+1)!}{1!}(x-x_0)\text{.}$$
	Logo,
	$$E_n^{(k)}(x)=f^{(k+1)}(\xi)\cdot (x-x_0)$$
	
	É possível concluir, portanto, que a $n$-ésima derivada de $P_n(x)$ é uma boa aproximação da $n$-ésima derivada da função $f(x)$, de modo que $E_n^{(n)}=o(x-x_0)$.
\end{document}